\subsection{Modified coefficient-domain protocol with shared PCS coordinates}
\label{sec:coeff-domain-modified}

We describe a variant of the coefficient-domain protocol (Protocol~A
in~\cref{sec:coeff-vs-crt}) that recovers the shared-coordinate PCS
advantage of the CRT-domain protocol, at the cost of $3\nu$ additional
sumcheck elements.  The key idea is to \emph{defer} the oracle queries
on $\tilde{t}$ and $\tilde{s}$ and batch their verification into
Sumcheck~2, so that all three final oracle evaluation points share a
common last-$\nu$ challenge~$\gamma$.

\parhead{Setup}
As before: $\ring_p = \ZZ_p[X]/(X^n+1)$, $n = 2^\nu$,
$A \in \ring_p^{N \times M}$, $s \in \ring_p^M$, $t \in \ring_p^N$
with $t = As$.  The prover commits to coefficient MLEs
$\tilde{A}(\bm{I},\bm{J},\bm{K})$,
$\tilde{s}(\bm{J},\bm{K})$,
$\tilde{t}(\bm{I},\bm{K})$.
Let $d_f(\cdot\,;\alpha',\beta')$ denote the $\nu$-variate shift
function arising from the negacyclic circulant decomposition,
evaluable by the verifier in $O(\nu)$ field operations.

\begin{figure}[H]
\caption{Protocol $\Pi_{\mathsf{RingMul}}^{\mathsf{coeff*}}$: modified
coefficient-domain PIOP with shared PCS coordinates.}
\label{prot:coeff-modified}
\begin{algorithm}[H]
\textbf{Input:} Oracles
$\oracle{\tilde{A}} \in \field_p^{\multilin}[\bm{I},\bm{J},\bm{K}]$,
$\oracle{\tilde{s}} \in \field_p^{\multilin}[\bm{J},\bm{K}]$,
$\oracle{\tilde{t}} \in \field_p^{\multilin}[\bm{I},\bm{K}]$.

\begin{algorithmic}[1]

\STATE The verifier samples
$\alpha \randpick \extensionfield^\tau$,
$\alpha' \randpick \extensionfield^\nu$.

\STATE \label{step:coeff-et}
The prover sends $e_t$ claimed to equal $\tilde{t}(\alpha, \alpha')$.

\STATE \label{step:coeff-sc1}
\textbf{Sumcheck~1} (degree-$2$, $\mu + \nu$ rounds).
$\prover$ and $\verifier$ run a sumcheck for the claim
\[
    e_t \;=\; \sum_{(j,k) \in \bits^{\mu+\nu}}
    V(j,k) \cdot \tilde{s}(j,k),
\]
where $V(j,k) = \sum_\ell \meq(\ell, \alpha')
\cdot [Z_{-1}(\tilde{A}(\alpha, j, \cdot))]_{\ell,\, k}$
is the negacyclic-circulant virtual polynomial.
Let $(\beta, \beta') \in \extensionfield^{\mu+\nu}$ denote the
folding challenges, reducing to the claim
$\sigma' = V(\beta,\beta') \cdot \tilde{s}(\beta,\beta')$.

\STATE \label{step:coeff-eV}
The prover sends $e_V$ claimed to equal $V(\beta,\beta')$.
The verifier computes $e_s := \sigma' / e_V$
(rejecting if $e_V = 0$ and $\sigma' \neq 0$).

\STATE \label{step:coeff-rho}
The verifier samples batching challenge
$\rho \randpick \extensionfield$.

\STATE \label{step:coeff-sc2}
\textbf{Sumcheck~2} (degree-$2$, $\nu$ rounds).
$\prover$ and $\verifier$ run a sumcheck for the batched claim
\begin{equation}\label{eq:coeff-batched-sc2}
    e_V + \rho\, e_t + \rho^2 e_s
    \;=\;
    \sum_{\ell \in \bits^\nu} \Big[
        \tilde{A}(\alpha,\beta,\ell) \cdot d_f(\ell;\, \alpha',\beta')
        \;+\; \rho\, \meq(\ell,\alpha') \cdot \tilde{t}(\alpha,\ell)
        \;+\; \rho^2\, \meq(\ell,\beta') \cdot \tilde{s}(\beta,\ell)
    \Big].
\end{equation}
Each term in the summand is a product of two multilinears in~$\ell$
(degree~$2$).
Let $\gamma \in \extensionfield^\nu$ denote the folding challenges,
reducing to the claim
\[
    c \;=\;
    \tilde{A}(\alpha,\beta,\gamma) \cdot d_f(\gamma)
    + \rho\, \meq(\gamma,\alpha') \cdot \tilde{t}(\alpha,\gamma)
    + \rho^2\, \meq(\gamma,\beta') \cdot \tilde{s}(\beta,\gamma).
\]

\STATE \label{step:coeff-queries}
The verifier evaluates $d_f(\gamma)$,
$\meq(\gamma,\alpha')$, and $\meq(\gamma,\beta')$
in $O(\nu)$ each, then queries
\[
    e'_A \gets \oracle{\tilde{A}(\alpha,\beta,\gamma)}, \qquad
    e'_t \gets \oracle{\tilde{t}(\alpha,\gamma)}, \qquad
    e'_s \gets \oracle{\tilde{s}(\beta,\gamma)},
\]
and checks
$c = e'_A \cdot d_f(\gamma) + \rho\, \meq(\gamma,\alpha')\, e'_t
     + \rho^2\, \meq(\gamma,\beta')\, e'_s$.

\end{algorithmic}
\end{algorithm}
\end{figure}


\parhead{Shared coordinates}
All three oracle queries share $\gamma$ as their last $\nu$
coordinates:
\begin{center}
\begin{tabular}{lll}
    Oracle & Point & Dimension \\
    \hline
    $\tilde{t}$ & $(\alpha, \gamma)$ & $\tau + \nu$ \\
    $\tilde{s}$ & $(\beta, \gamma)$ & $\mu + \nu$ \\
    $\tilde{A}$ & $(\alpha, \beta, \gamma)$ & $\tau + \mu + \nu$
\end{tabular}
\end{center}
This is the same shared-suffix structure as the CRT-domain
protocol (\cref{prot:ringmul}), so the full PIOR consolidation
chain of \cref{sec:reducing-commitments} (SelectorMerge, MultiPoint,
DP24~packing, interleaved BaseFold) applies identically.

\parhead{Costs}
\begin{center}
\begin{tabular}{lccc}
    & Coeff.\ (original) & Coeff.\ (modified) & CRT \\
    \hline
    Sumcheck~1 & $3(\mu+\nu)$ & $3(\mu+\nu)$ & $3\mu$ \\
    Sumcheck~2 & $3\nu$ & $3\nu$ & $3\nu$ \\
    Prover claims & 1 & 2 & 2 \\
    \textbf{Total} & $3\mu + 6\nu + 1$ & $3\mu + 6\nu + 2$ & $3\mu + 3\nu + 2$ \\
    \hline
    Shared last-$\nu$ PCS & No & Yes & Yes \\
    PIOR chain & Separate opens & Full consolidation & Full consolidation
\end{tabular}
\end{center}
The modification adds one prover claim ($e_t$) and trades the three
unrelated PCS evaluation points of the original coefficient protocol
for a shared-$\gamma$ suffix, at no change in sumcheck round count or
degree.  The residual gap with the CRT protocol is exactly $3\nu$
sumcheck elements, arising entirely from Sumcheck~1 ranging over
$\mu + \nu$ variables (convolution entangles column and coefficient
indices) rather than $\mu$.

\begin{theorem}[Completeness]
If $t = As$ over $\ring_p$ and both parties follow the protocol,
then the verifier accepts.
\end{theorem}

\begin{proof}
Assume $t = As$.

\emph{Step~\ref{step:coeff-et}.}
The honest prover sets $e_t = \tilde{t}(\alpha,\alpha')$.

\emph{Step~\ref{step:coeff-sc1}.}
The negacyclic convolution gives
$t_i[\ell] = \sum_j [Z_{-1}(A_{ij})]_{\ell,k}\, s_j[k]$
for all $(i,\ell)$.  Passing to MLEs and evaluating at $(\alpha,\alpha')$:
\[
    \tilde{t}(\alpha,\alpha')
    = \sum_{(j,k) \in \bits^{\mu+\nu}} V(j,k) \cdot \tilde{s}(j,k),
\]
so the sumcheck claim $e_t = \text{RHS}$ holds.  By completeness
of degree-$2$ sumcheck, the honest round messages are accepted.

\emph{Step~\ref{step:coeff-eV}.}
The sumcheck reduces to $\sigma' = V(\beta,\beta') \cdot \tilde{s}(\beta,\beta')$.
The honest prover sets $e_V = V(\beta,\beta')$, so
$e_s = \sigma'/e_V = \tilde{s}(\beta,\beta')$.

\emph{Step~\ref{step:coeff-sc2}.}
We verify the three terms on the RHS of~\eqref{eq:coeff-batched-sc2}
separately.  The first:
\[
    \sum_{\ell \in \bits^\nu} \tilde{A}(\alpha,\beta,\ell) \cdot d_f(\ell)
    \;=\; V(\beta,\beta') \;=\; e_V,
\]
by the definition of $d_f$ (encoding the negacyclic shift at
challenges $\alpha',\beta'$).  The second:
\[
    \sum_{\ell \in \bits^\nu} \meq(\ell,\alpha') \cdot \tilde{t}(\alpha,\ell)
    \;=\; \tilde{t}(\alpha,\alpha'),
\]
since $\sum_\ell \meq(\ell,\alpha')\, f(\ell) = f(\alpha')$ for any
multilinear~$f$.  Hence this contributes $\rho \cdot e_t$.
The third:
\[
    \sum_{\ell \in \bits^\nu} \meq(\ell,\beta') \cdot \tilde{s}(\beta,\ell)
    \;=\; \tilde{s}(\beta,\beta') \;=\; e_s,
\]
contributing $\rho^2 \cdot e_s$.  The total is
$e_V + \rho\, e_t + \rho^2 e_s$, matching the LHS.

\emph{Step~\ref{step:coeff-queries}.}
By completeness of Sumcheck~2, the reduced claim at $\gamma$ holds.
The oracle queries return the honest evaluations, and the verifier's
check passes.
\end{proof}


\begin{theorem}[Soundness]
Protocol $\Pi_{\mathsf{RingMul}}^{\mathsf{coeff*}}$ has
round-by-round knowledge soundness error at most
\[
    \epsilon_{\mathsf{rbr}}
    \;\leq\;
    \underbrace{\frac{\tau + 2\nu}{|\extensionfield|}}_{\text{Schwartz--Zippel}}
    \;+\; \underbrace{\frac{2(\mu+\nu)}{|\extensionfield|}}_{\text{Sumcheck 1}}
    \;+\; \underbrace{\frac{2}{|\extensionfield|}}_{\text{batching}}
    \;+\; \underbrace{\frac{2\nu}{|\extensionfield|}}_{\text{Sumcheck 2}}
    \;=\;
    \frac{\tau + 2\mu + 6\nu + 2}{|\extensionfield|}.
\]
For comparison, the CRT protocol (\cref{thm:ringmul-rbr}) achieves
$(\tau + 2\mu + 4\nu + 1)/|\extensionfield|$; the difference of
$2\nu + 1$ arises from the $\nu$ extra Sumcheck~1 rounds and the
degree-$2$ (versus degree-$1$) batching polynomial.
\end{theorem}

\begin{proof}[Proof sketch]
The protocol has four phases contributing to the soundness error.

\emph{(i) Schwartz--Zippel for the main relation.}
Define $R(i,\ell) := t_i[\ell] - \sum_j \sum_k
[Z_{-1}(A_{ij})]_{\ell,k}\, s_j[k]$ for $(i,\ell) \in [N] \times [n]$,
and let
$H(\alpha,\alpha') := e_t - \sum_{(j,k)} V(j,k)\, \tilde{s}(j,k)$.
On Boolean inputs $H$ agrees with $\tilde{R}$, so $R \not\equiv 0$
implies $H \not\equiv 0$.  Since $H$ has individual degree~$1$ in
the $\tau$ variables of~$\alpha$ and individual degree~$2$ in
the $\nu$ variables of~$\alpha'$ (from the product structure of the
circulant), Schwartz--Zippel gives
$\Pr[H(\alpha,\alpha') = 0] \leq (\tau + 2\nu)/|\extensionfield|$.

\emph{(ii) Sumcheck~1 RBR soundness.}
Degree-$2$ sumcheck over $\mu + \nu$ rounds contributes
$2(\mu+\nu)/|\extensionfield|$.

\emph{(iii) Batching.}
After Sumcheck~1, the prover commits to $e_V$ and the verifier
derives $e_s = \sigma'/e_V$.  Define the errors
$\Delta_V = e_V - V(\beta,\beta')$,
$\Delta_t = e_t - \tilde{t}(\alpha,\alpha')$,
$\Delta_s = e_s - \tilde{s}(\beta,\beta')$.
The batched error is $\Delta_V + \rho\,\Delta_t + \rho^2\Delta_s$,
a degree-$2$ polynomial in the random $\rho$ (sampled after $e_V$
is fixed).  If any $\Delta \neq 0$, this vanishes with probability
$\leq 2/|\extensionfield|$.

Since $e_s = \sigma'/e_V$ and Sumcheck~1 fixes $\sigma'$,
an incorrect $e_V \neq V(\beta,\beta')$ generically forces
$e_s \neq \tilde{s}(\beta,\beta')$, so at least one of $\Delta_V$,
$\Delta_s$ is nonzero.

\emph{(iv) Sumcheck~2 RBR soundness.}
Degree-$2$ sumcheck over $\nu$ rounds contributes
$2\nu/|\extensionfield|$.

Summing all terms gives the stated bound.
\end{proof}


\parhead{Integration with SpeedySpartan}
In a SpeedySpartan-style outer protocol, the inner sumcheck
(our Sumcheck~1) produces challenges $(\beta,\beta')$ that also
determine the evaluation points for the wiring matrices
$\tilde{W}_w(r,\beta,\beta')$.  These claims can be absorbed into
Sumcheck~2 at no increase in degree by adding terms
$\rho_w\, \meq(\ell,\beta') \cdot \tilde{W}_w(r,\beta,\ell)$
to~\eqref{eq:coeff-batched-sc2}.  After folding with~$\gamma$,
the wiring-matrix queries land at $(r,\beta,\gamma)$---sharing
$\gamma$ with the ring-IOP claims---so the full PIOR consolidation
chain applies to all oracles simultaneously.

The gate selector $\tilde{q}_M(r)$ is lower-dimensional and
independent of the coefficient-index variables;
padding it to $(r,\beta,\gamma)$ (constant in $(\bm{J},\bm{L})$)
makes it compatible with any shared suffix.

\emph{SelectorMerge alignment.}
After Sumcheck~2 with challenge~$\gamma$, the evaluation claims are:
\begin{center}
\begin{tabular}{llll}
    Oracle & Point & Dim. & Last $\mu{+}\nu$ coords \\
    \hline
    $\tilde{W}_w$ ($w \in \{a,b,c\}$) & $(r, \beta, \gamma)$ & $m + \mu + \nu$ & $(\beta, \gamma)$ \\
    $\tilde{q}_{M,\mathrm{pad}}$ & $(r, \beta, \gamma)$ & $m + \mu + \nu$ & $(\beta, \gamma)$ \\
    $\tilde{A}$ & $(\alpha, \beta, \gamma)$ & $\tau + \mu + \nu$ & $(\beta, \gamma)$ \\
    $\tilde{t}_{\mathrm{pad}}$ & $(\alpha, \beta, \gamma)$ & $\tau + \mu + \nu$ & $(\beta, \gamma)$ \\
    $\tilde{s}$ & $(\beta, \gamma)$ & $\mu + \nu$ & $(\beta, \gamma)$
\end{tabular}
\end{center}
Here $\tilde{t}_{\mathrm{pad}}(x,j,\ell) := \tilde{t}(x,\ell)$
is padded to be constant in~$\bm{J}$, so
$\tilde{t}_{\mathrm{pad}}(\alpha,\beta,\gamma) = \tilde{t}(\alpha,\gamma)$.
All seven claims share $(\beta,\gamma)$ as their last $\mu + \nu$
coordinates.  The indexed oracles split into two families:
\begin{enumerate}[nosep]
    \item $\tilde{W}_a, \tilde{W}_b, \tilde{W}_c, \tilde{q}_{M,\mathrm{pad}}$
          at data point $(r, \beta, \gamma)$---SelectorMerge with $2$
          selector bits covers $\bits^2$ exactly, so MultiPoint
          costs $0$ prover messages.
    \item $\tilde{A}, \tilde{t}_{\mathrm{pad}}$
          at data point $(\alpha, \beta, \gamma)$---SelectorMerge
          with $1$ selector bit covers $\bits^1$ exactly,
          MultiPoint costs $0$ prover messages.
\end{enumerate}
After MultiPoint, there are $3$ claims: $h_1(r'_1, r'_2, r, \beta, \gamma)$,
$h_2(r'_3, \alpha, \beta, \gamma)$, and $\tilde{s}(\beta, \gamma)$.
These share $(\beta,\gamma)$ as a suffix but differ in prefix
coordinates $(r'_1, r'_2, r)$ vs.\ $(r'_3, \alpha)$.
They cannot be further merged by SelectorMerge (different prefix
dimensions and values), so they proceed as three independent claims
into DP24~packing and BaseFold.

However, interleaved BaseFold still benefits: once each polynomial
is solo-folded down to $\mu + \nu$ variables, all three share the
same folding challenges for the remaining $\mu + \nu$ rounds,
and a single interleaved Merkle tree covers all three simultaneously.

\emph{Comparison with the CRT variant.}
The CRT SpeedySpartan instance (\cref{sec:reducing-commitments})
produces $2$ claims after PIOR consolidation ($h_1$ and $\tilde{s}$),
because $\tilde{A}$ and $\tilde{t}$ are handled by the CRT
consistency sumcheck rather than as separate oracle claims.
The coefficient variant produces $3$ claims, adding one
indexed-oracle opening ($h_2$) to the BaseFold phase.
The concrete cost of this extra opening is $1 + \tau$ solo-fold
rounds for $h_2$ (dimension $1 + \tau + \mu + \nu$, shared
dimension $\mu + \nu$) before joining the interleaved phase.
For SWIFFT ($\tau = 4$): $5$ extra folding rounds---each
contributing one Merkle commitment and $q$ authentication
paths---versus the CRT variant.

\parhead{Comparison with CRT protocol}
The PCS costs are now identical.  The irreducible difference is
$3\nu$ sumcheck extension-field elements, arising from Sumcheck~1's
$\mu + \nu$ variables versus the CRT protocol's $\mu$.
For SWIFFT parameters ($\mu = 4$, $\nu = 6$):
$3\nu = 18$ extension-field elements $\approx 90$ bytes---negligible
relative to hash-based PCS proof sizes of tens of kilobytes.
The CRT-domain protocol's advantage in this regime is therefore
structural simplicity (one sumcheck architecture, no shift function
$d_f$) rather than concrete proof size.